\documentclass[10pt,a4paper]{amsart}
\usepackage[margin=19mm]{geometry}
\usepackage{amsmath,amssymb,mathtools}
\usepackage[scr=rsfs,cal=euler]{mathalfa}
\usepackage{xfrac}
\usepackage[unicode,hidelinks,bookmarks=false]{hyperref}
\newcommand{\ie}{i.e.\ }
\newcommand{\cf}{cf.\ }
\newcommand{\resp}{respectively\ }
\newcommand{\Subsection}{\subsection}
\providecommand{\nobreakdash}{\nobreak\hbox{-}}
\setlength{\emergencystretch}{2em}
\multlinegap0pt
\usepackage{amssymb}
\usepackage{float}
\usepackage{cancel}
\usepackage{algorithm}\usepackage{algpseudocode}
\usepackage{tikz}
\newtheorem{cor}{Corollary}
\newtheorem{lem}{Lemma}
\title{Theory and Computation of Discrete Sturm--Liouville Problems on Non-Uniform Grids via the Pr\"ufer Transformation}
\author{Bailyn Hall, Kimsear Lor}
\date{Source: arXiv:2607.15300v1 (CC BY 4.0)}
\begin{document}
\maketitle

\noindent\emph{Source and licence.} Theory and Computation of Discrete Sturm--Liouville Problems on Non-Uniform Grids via the Pr\"ufer Transformation. Version arXiv:2607.15300v1 (CC BY 4.0), distributed by arXiv under the Creative Commons Attribution 4.0 International licence, http://creativecommons.org/licenses/by/4.0/, as recorded in the arXiv licence metadata for this exact version and shown on the abstract page https://arxiv.org/abs/2607.15300v1. Full licence notice: \texttt{licenses/arxiv-2607.15300-CC-BY-4.0.txt}.\par\smallskip

\noindent\textit{Editorial scope.} Counted unit: the article's lem lem:opL0 (source0.tex lines 217--224). The article's complete original proof of it is delivered with it (source0.tex lines 225--237, one \texttt{proof} environment of 13 lines). No mathematics is written, repaired or re-derived: the statement and the proof are the article's own lines, in the article's own notation, and no retained line is substituted or re-flowed.  Retained uncounted as the source-local context the proof needs: lines 85--88; lines 170--176; lines 191--194; lines 212--215.  Nothing was omitted from any retained span, so the package carries no omission record.  No retained line contains a citation, so the wrapper carries no bibliography and no bibliography entry.  No retained line contains a figure environment, an image-inclusion command or drawing code, so no image is delivered and the package declares no asset dependency.  Editorial formatting only, in addition: an AMS \texttt{amsart} preamble that restates the article's own declarations and macros used by the retained lines, namely the environments \texttt{cor}, \texttt{lem} on separate counters, together with the standard packages those lines need.  The retained environments print in this document as Corollary 1, Lemma 1 in order of appearance, whereas the article prints the same items with the numbering of its own full document.  The complete native source of the article as distributed in the arXiv e-print for this exact version is bundled unchanged as \texttt{sources/205376/source0.tex}.
\begin{equation}
\label{def:2.2}
\Delta_h f_k = \frac{f_{k+1} - f_k}{t_{k+1} - t_k} = \frac{f_{k+1} - f_k}{h_k},    
\end{equation}

\begin{cor}
\label{cor:summation}
Let $\{a_k\}$ and $\{b_k\}$ be sequences of real numbers, with $h_k$ defined as in \eqref{def:2.2}. If $b_0 = b_{n+1} = 0$, then
\[
\sum_{k=0}^{n} h_k\, b_{k+1}\, \Delta_h a_k = -\sum_{k=0}^{n} h_k\, a_k\, \Delta_h b_k.
\]
\end{cor}

\begin{equation}
\label{eq:inner-product}
\left\langle u,v\right\rangle := \sum_{k=0}^{n} h_k\, u_{k+1} \overline{v_{k+1}}.
\end{equation}

\begin{equation}
    \label{def:op}
    (\mathcal{L}_0 x)_{k+1} = \Delta_h(r_k \Delta_h x_k).
\end{equation}

\begin{lem}
\label{lem:opL0}
The principal operator $\mathcal{L}_0$ defined in \eqref{def:op} is self-adjoint with respect to the inner product \eqref{eq:inner-product}; that is,
\[
\langle \mathcal{L}_0 x, y \rangle = \langle x, \mathcal{L}_0 y \rangle
\]
for all sequences $x, y$ satisfying the Dirichlet boundary conditions.
\end{lem}

\begin{proof}
Observe that $\langle \mathcal{L}_0x,y \rangle = \displaystyle\sum_{k=0}^{n}h_k\Delta_h(r_k \Delta_h x_{k})\overline{y_{k+1}}$ and $\langle x,\mathcal{L}_0y \rangle = \displaystyle\sum_{k=0}^{n}h_k\Delta_h(r_k \Delta_h y_{k})\overline{x_{k+1}}$. Let 
\begin{align*}
a_k &= r_k\Delta_h x_k,  \qquad \widetilde{a}_k = r_k\Delta_h y_k,\\
b_k &= \overline{y_k}, \qquad \widetilde{b}_k=\overline{x_k}.
\end{align*}
By employing Corollary \ref{cor:summation}, we obtain
\begin{align*}
    \displaystyle\sum_{k=0}^{n}h_k\Delta_h(r_k\Delta_hx_k)\overline{y_{k+1}} &= -\displaystyle\sum_{k=0}^{n}h_k(r_k\Delta_hx_k)\Delta_h\overline{y_k},\\
     \displaystyle\sum_{k=0}^{n}h_k\Delta_h(r_k\Delta_hy_k)\overline{x_{k+1}} &= -\displaystyle\sum_{k=0}^{n}h_k(r_k\Delta_hy_k)\Delta_h\overline{x_k},
\end{align*}
which imply $\displaystyle\sum_{k=0}^{n}h_k\Delta_h(r_k \Delta_h x_{k})\overline{y_{k+1}}-\displaystyle\sum_{k=0}^{n}h_k\Delta_h(r_k \Delta_h y_{k})\overline{x_{k+1}}=0$, using that $r_k$ is real. Therefore, we conclude $\langle \mathcal{L}_0x,y \rangle= \langle x,\mathcal{L}_0y \rangle$.
\end{proof}

\end{document}
